% Abstract content - included via \input in main.tex's abstract environment
One of the questions which software product teams are likely to face in this era of high technological change is how and what to do with such a multitude of various kinds of features prior to being able to allocate resources based on prospective information regarding the demand in a product design process.

The focal point of the proposed study is Notion, which is a tremendously popular SaaS (Software as a Service) productive platform, as data analytics models owned by the proprietors cannot be shared publicly, and, instead, the data on search interest in Google Trends, app store review corpora, and user-feedback patterns can be publicly available to generate a specific demand signal of a particular feature of a product.

The proposed forecasting system would employ an ensemble technique of sampling that incorporates different methods of analysis, which include Exponential Smoothing, ARIMA and linear regression models. The different approaches would thus enable the different concerned parties to have a picture of the larger picture.

As the empirical evidence indicates, it is evident that it is possible to recognise the recognisable trends in the analysed attributes, including AI-driven features, template features, and automation processes. The trends in the development of AI-related factors include long-term growth curves, templates possess stable bases of demand, with predictable seasonal variation, and automation capability, and the automation features have intermediate growth rates with regard to the workflow optimization trends in the industry.

The core worth on the work is the clarification that one can only identify meaningful demand forecasting through pure utilisation of publicly accessible data in determining the prioritisation of software features. The explanation is that companies with limited analytical ability can only reduce the extent to which development resources are misplaced and also improve the alignment between product roadmap and user demands.

\textbf{Index Terms}---SaaS product management, demand forecasting, time series analysis, feature prioritisation, Google Trends, sentiment analysis, ARIMA, exponential smoothing.
